\section{Related work}\label{sec:related-work}

We position the proposed method against six families of hyperparameter
optimization (HPO) algorithms, plus the OpenML benchmarking ecosystem
that defines our evaluation protocol. The selected baselines span
single-objective HPO under a fixed evaluation budget, multi-fidelity
optimization, evolutionary search, and multi-objective HPO; AutoML
systems are referenced as context rather than as direct baselines.

\subsection{Classical HPO baselines}\label{sec:rw-classical}
The dominant single-objective baselines for tabular HPO are random
search \citep{bergstra2012random}, Gaussian-process Bayesian
optimization (GP-BO) \citep{snoek2012bayesian,shahriari2016bayesian,
frazier2018tutorial}, and the tree-structured Parzen estimator (TPE)
\citep{bergstra2011tpe}, the latter as implemented by Hyperopt and
exposed in the form of the Optuna \texttt{TPESampler}
\citep{akiba2019optuna}. \citet{probst2019tunability} report that for
gradient-boosting decision trees (GBDT) most of the variance in
generalization is captured by a small set of tuned hyperparameters,
which makes random search a non-trivial control even at modest
budgets.

\subsection{Sequential model-based algorithm
configuration}\label{sec:rw-smac}
The Sequential Model-Based Algorithm Configuration (SMAC) family
\citep{hutter2011smac,lindauer2022smac3} replaces the GP surrogate
with a random forest, which natively handles mixed-type and
conditional hyperparameter spaces and tolerates noisy evaluations via
intensification and racing. SMAC3 is the de facto state of the art
for general-purpose HPO of tabular ML pipelines, and powers the
classical Auto-sklearn \citep{feurer2015autosklearn} backend.

\subsection{Multi-fidelity HPO}\label{sec:rw-multifidelity}
Multi-fidelity methods reuse cheap, low-fidelity evaluations
(e.g., few epochs, few CV folds, sub-sampled training set) to
prune unpromising configurations early. Successive Halving and
Hyperband \citep{li2017hyperband} formalize this as a bandit problem;
ASHA \citep{li2020asha} parallelizes the schedule for asynchronous
workers. BOHB \citep{falkner2018bohb} couples Hyperband's bracket
schedule with TPE-style model-based sampling, and DEHB
\citep{awad2021dehb} replaces TPE with differential evolution
\citep{storn1997differential} for better high-dimensional behavior.
For GBDT models, the most natural fidelity dimension is the number of
boosting iterations (\texttt{n\_estimators} for XGBoost / LightGBM /
CatBoost), which lets multi-fidelity methods reuse partially trained
ensembles cheaply.

\subsection{Evolutionary and derivative-free
HPO}\label{sec:rw-evolutionary}
Differential evolution \citep{storn1997differential} and
covariance-matrix adaptation evolution strategy (CMA-ES)
\citep{hansen2001cmaes} are derivative-free baselines for continuous
HPO; they are particularly relevant when the objective surface is
non-smooth or when the budget allows a generation-based search rather
than the strictly sequential single-incumbent dynamics of GP-BO. The
Nevergrad platform \citep{rapin2018nevergrad} is referenced as a
unified implementation harness for evolutionary and derivative-free
optimizers; the present article cites it for context rather than
reporting Nevergrad-only baselines.

\subsection{Multi-objective HPO}\label{sec:rw-moo}
\citet{morales2022survey} survey multi-objective HPO and group the
literature into evolutionary, scalarization-based, and decomposition
approaches. NSGA-II \citep{deb2002nsga2} and its many-objective
extension NSGA-III \citep{deb2014nsga3} are the canonical
evolutionary baselines (the latter is relevant for $q\!>\!2$
objectives via reference points, paralleling the simplex-lattice
weight grid we use for NBI). The multi-objective TPE variant MOTPE
\citep{ozaki2020motpe} is the natural multi-objective counterpart of
TPE-Hyperopt. ParEGO \citep{knowles2006parego} represents
scalarization-based BO with rotating weight vectors. Within the
DoE / RSM tradition, multi-objective problems are commonly attacked
with Normal Boundary Intersection \citep{das1998nbi}, which generates
evenly spaced points on the Pareto front of nonlinear MOO problems
and is well-defined for arbitrary $q$ objectives. The pymoo library
\citep{blank2020pymoo} provides reference Python implementations of
NSGA-II / NSGA-III and is the natural implementation candidate for the
evolutionary baselines in the comparative protocol.

\subsection{Multivariate NBI with factor
scores}\label{sec:rw-pereira}
The closest methodological reference for the article-track is
\citet{pereira2025eaai}, who replace the original responses with
Varimax-rotated factor scores ($\VRF$), wrap each $\VRF$ in an FMSE
target-seeking loss, and solve the resulting NBI sub-problem under a
spherical region constraint. They further introduce a
mixture-design-based post-optimization that fits Scheff\'{e}
canonical polynomials over the simplex of NBI weights and refines
$\bm w^{*}$ by minimizing a generalized distance to utopia subject to
an elliptical interior constraint and an entropy preference.

\subsection{AutoML systems as context}\label{sec:rw-automl}
Auto-sklearn \citep{feurer2015autosklearn}, FLAML
\citep{wang2021flaml} and AutoGluon-Tabular
\citep{erickson2020autogluon} bundle HPO with model selection,
ensembling, and resource control. They are cited as context: the
proposed method targets the HPO sub-problem of a fixed GBDT
algorithm, not full pipeline AutoML. FLAML is the most natural
practical baseline of the three (single-algorithm tuning with a CPU
budget) and is listed in the comparative protocol as an optional
addition; Auto-sklearn and AutoGluon are excluded from the headline
comparison because their search spaces are not directly aligned with
the per-algorithm hyperparameter list of XGBoost / LightGBM /
CatBoost.

\subsection{Benchmark suites}\label{sec:rw-benchmarks}
We use the OpenML-CC18 curated classification benchmark suite
\citep{bischl2021openmlbenchmark} as the primary evaluation panel.
OpenML-CC18 (\texttt{suite\_id = 99}) collects 72 standardized
classification tasks selected for license clarity, dataset diversity,
and well-defined train/test splits, and is widely used as a fair
common ground for HPO and AutoML evaluation. The OpenML platform
\citep{vanschoren2014openml} provides task identifiers, default folds,
and metadata that make task-based (rather than dataset-based)
benchmarks reproducible across implementations. Cross-method
statistical comparison follows \citet{demsar2006statistical}.

\subsection{This work in context}\label{sec:rw-positioning}
Relative to \citep{pereira2025eaai}, we (a) port the framework to HPO
of GBDT models with explicit hyperparameter box constraints rather
than a hyperspherical region; (b) keep the FMSE wrapping for $\VRF$
objectives but support a parallel direct-objective mode for raw ML
metrics with explicit directions; (c) make the NBI sub-problem the
\emph{actual} solver in the implementation rather than only a
mathematical reference. Relative to
\citep{ribeiro2026dissertation}, we replace the legacy weighted-sum
scalarization with the true $N$-objective NBI, generalize the
2-VRF / 3-VRF collapse to a configurable factor model
(\textsc{auto}/\textsc{fixed}/\textsc{manual}/\textsc{none}), make
RSM design-aware (separating process-quadratic from Scheff\'{e}
mixture models), and add the conditional MBPA stage. The legacy
weighted-sum solver is preserved as an explicit ablation baseline and
is never referred to as NBI in the article-track code or text. The
formal comparative protocol against the baselines listed above is
specified in \texttt{docs/COMPARATIVE\_PROTOCOL.md} and the
machine-readable method matrix at
\texttt{benchmarks/doctoral/openml\_cc18/method\_matrix.csv}.
